% !TEX root = ../../main.tex
% C8.5-C8.7 - Performance and camera stream handling (translated). Figures: report-data-guide.md C8.5-C8.7 (updated 2026-09-29),
% with the mandatory measurement conditions.
\section{Performance and Camera Stream Handling}
\label{sec:8.5}
\label{sec:8.6}
\label{sec:8.7}

\paragraph{Sampling interval.} Table~\ref{tab:sampling} compares the two sampling presets on the same 10-second video (600 frames) with the current configuration, each preset run 3 times, interleaved, on the target machine. The first run of $N = 10$ had to load the models (taking 178.5 seconds), so the time of $N = 10$ is the average of the other two runs. Other applications on the machine were not closed during the measurement, so the time varied by about $\pm 12$ seconds between runs. CPU and memory were sampled every 0.25 seconds, so they are lower bounds of the actual values.

\begin{table}[H]
\centering
\caption{Comparison of the two sampling intervals on a 10-second video (confidence threshold 0.1)}
\label{tab:sampling}
\begin{tabularx}{\textwidth}{|L|>{\centering\arraybackslash}p{3.4cm}|>{\centering\arraybackslash}p{3.4cm}|}
\hline
\textbf{Metric} & \textbf{$N = 10$} & \textbf{$N = 20$} \\
\hline
Processing time with the models loaded & 149.5 s & 132.6 s \\
\hline
CPU time (main process and model processes) & about 212 s & about 160 s \\
\hline
Peak memory of the main process & 695 MiB & 576 MiB \\
\hline
Peak memory of the model processes & 4,284 MiB & 4,267 MiB \\
\hline
Tracks & 44 & 36 \\
\hline
Short tracks (under 2 seconds) & 22 & 24 \\
\hline
Size of the representative images & 10.2 MB & 6.9 MB \\
\hline
\end{tabularx}
\end{table}

With $N = 20$, CPU time drops by about a quarter and the image size by about a third, while the number of short tracks, a sign of broken tracks, hardly changes. Processing time drops only about 9\%, because most of it is spent in the image encoder, about 1.85 seconds per track against 0.16 to 0.23 seconds per sampled frame for the Detector; it therefore depends mainly on the number of people in the video. The benchmark video is more crowded than the demonstration data: 36 tracks in 10 seconds, against 752 tracks in 420 seconds of video from the seven cameras in the 60-120 second segment (Section~\ref{sec:7.2}). The ratio of processing time to video length is therefore higher here than the roughly sevenfold slowdown measured when indexing the demonstration data. Model process memory barely differs, since it mostly holds weights. As a CPU-only machine is already far slower than real time, $N = 20$ stays the default (Section~\ref{sec:3.4}).

\paragraph{RTSP stream handling.} The ingestion of simulated RTSP streams was verified in three situations, with a confidence threshold of 0.25, so the track counts are only indicative: a normal session processed 300 frames, produced 29 tracks and a text search returned 8 results, in 136 seconds; when the streaming source was cut for 5 seconds mid-session, the background worker reconnected once and continued, completing 600 frames with 38 tracks in 140 seconds; when the Administrator disabled AI processing mid-session, the job became \emph{Cancelled} and left no incomplete track. With the current configuration, in the rehearsals, a newly added RTSP camera was \emph{Online} and the background worker created a processing session 13 to 24 seconds after AI processing was enabled.

\paragraph{Search latency.} Latency was measured on the client side by calling the search API with $k = 8$ on the 1,523-track demonstration data, 30 calls per form after one warm-up call, interleaved and 2.5 seconds apart, with the background worker idle. The results are in Table~\ref{tab:latency}; no call was rejected by the rate limit.

\begin{table}[H]
\centering
\caption{Search latency through the API (30 calls per form, $k = 8$)}
\label{tab:latency}
\begin{tabularx}{\textwidth}{|L|c|c|c|}
\hline
\textbf{Form} & \textbf{Median} & \textbf{95th percentile} & \textbf{Maximum} \\
\hline
By image & 1,269 ms & 1,976 ms & 2,289 ms \\
\hline
By text & 106 ms & 142 ms & 151 ms \\
\hline
By attributes & 103 ms & 136 ms & 136 ms \\
\hline
\end{tabularx}
\end{table}

Search by image is over ten times slower because the query image goes through RaSa's image encoder on the CPU, while a description only needs the text encoder. Through the interface in the rehearsals, the first search after start-up took about 31 seconds, mostly to load the encoder. Later searches took about 8 seconds, as they were measured while the worker ran an RTSP session and include the rendering time. Without warm-up and with a busy worker, the first search took about 97 seconds, which is why the deployment procedure includes a warm-up search (Section~\ref{sec:7.1}).

\paragraph{Storage and memory.} Each track takes about 358~KB on average for its full JPEG frame and 1~KB for its 256-dimensional vector; the relational database is about 24~MB. Memory was measured for each component with the whole system running as in the demonstration, taking the maximum in each state. When idle and after the search warm-up, the services in Docker use 515~MiB (828~MiB including the Docker virtual machine), the application server 1,978~MiB (153~MiB before loading the query encoder), the background worker 338~MiB, the front-end development server 318~MiB and the FFmpeg processes streaming the cameras 286~MiB, about 3.7~GiB in total. While searching, the application server rises to 2,037~MiB. While processing a job, the analysis pipeline uses up to 4,772~MiB more, bringing the system's total to about 8.0~GiB.

The memory figures come with four conditions. First, to avoid writing into the demonstration data, the processing state was measured with a tool that runs exactly the system's analysis pipeline on a 10-second video in place of the background worker. Second, the peak of 4,772~MiB adds the peaks of the main process and the model processes, which may not occur at the same time, so it is an upper bound. Third, on Windows, the measured memory is the part currently in RAM, which the operating system may trim when RAM runs low. Fourth, other applications were open during the measurement, so the RAM used by the whole machine rose from 13.2~GiB when idle to 13.9~GiB while searching and 15.2~GiB while processing, out of 15.7~GiB. The main reason is that the server and the pipeline load their own models, about 1.8~GiB for the query encoder and 4.2~GiB for the Detector and image encoder, so unneeded applications should be closed during a demonstration.
